We classify regular full-dimensional stochastic containment among binary
standard semi-directed strongly tree-child level-2 phylogenetic networks
under the Kimura three-parameter (K3P) model.  On the principal positive
Fourier domain, a directed containment germ exists if and only if the labelled
reduced trees of blobs agree and corresponding complete factors are either
labelled-isomorphic or ordinarily triangle-redirected, with coherent boundary
transports.  The same condition is equivalent to a common
full-dimensional regular germ, and it remains exact in strict continuous
time.  Thus no proper one-sided containment occurs in the strong class, and
the semi-directed topology is generically identifiable and, outside the
same proper exceptional set, exactly reconstructible modulo ordinary
triangle redirection.

K3P has a distinct local geometry.  The three ordinary triangle orientations
have generic normalized rank \(14\), share the same irreducible eight-term
quartic hypersurface \(H_{14}\subset\mathbb A^{15}\), and meet in a common
strict continuous-time smooth rank-\(14\) germ; they do not fill an
ambient-open three-leaf germ.  Necessity combines three-sector bridge fibres,
physical marginal submersions, a noncircular bridge-split transfer theorem,
and an exact finite classification.  The cut argument deliberately does not
assert a universal pointwise K3P rank criterion: true cuts have rank at most
four pointwise, noncuts are generically detected, and equality of cut sets is
proved under source-relative containment in the strong class by 204
pointwise one-active obstructions.  The remaining bounded residue consists
of fourteen four-port orbits---nine polynomially separated and five
directed-rank separated---plus two separately separated sink swaps.

The strong hypothesis is sharp.  For every \(n\ge3\), two weakly but not
strongly tree-child networks have strict continuous-time K3P images sharing
a common full-dimensional regular germ of dimension \(6n-3\).  The base
rank-\(15\) common point is certified by exact rational interval arithmetic
and a Krawczyk inclusion.  Exact certificates, separately implemented
replays, and fail-closed mutation tests accompany every finite and interval
claim.
